% !TeX root = ../Book5.tex
\begin{figure}[htbp]
\centering
\begin{tikzpicture}[x=1cm,y=1cm,every node/.style={font=\small},line width=.6pt]
  \begin{scope}
    \node at (2.5,1.65) {积分变量的 \(t\)-平面};
    \draw[->] (0,-1.15)--(0,1.25) node[above] {\(\operatorname{Im}t\)};
    \draw[->] (-.25,0)--(5.25,0) node[right] {\(\operatorname{Re}t\)};
    \node[below left] at (0,0) {\(0\)};
    \draw[blue!70!black,->] (0,.04)--(1.8,.04);
    \draw[blue!70!black] (1.8,0) arc[start angle=180,end angle=0,radius=.5];
    \draw[blue!70!black,->] (2.8,.04)--(5.05,.04);
    \draw[red!70!black,->] (0,-.04)--(1.8,-.04);
    \draw[red!70!black] (1.8,0) arc[start angle=180,end angle=360,radius=.5];
    \draw[red!70!black,->] (2.8,-.04)--(5.05,-.04);
    \draw[blue!70!black,->] (2.3,.5)--(2.38,.493);
    \draw[red!70!black,->] (2.3,-.5)--(2.38,-.493);
    \fill (2.3,0) circle (2pt);
    \node[below] at (2.3,-.08) {\(1\)};
    \node[blue!70!black] at (2.3,.83) {\(\gamma_+\)};
    \node[red!70!black] at (2.3,-.83) {\(\gamma_-\)};
    \node at (2.5,-1.55) {\(\gamma_+-\gamma_-\) 绕极点顺时针闭合};
  \end{scope}
  \begin{scope}[xshift=7.25cm]
    \node at (1,1.65) {解变量的 \(x\)-平面};
    \fill[blue!10] (0,0)--(58:1.7) arc[start angle=58,end angle=-58,radius=1.7]--cycle;
    \draw[->] (-.6,0)--(2.7,0) node[right] {\(\operatorname{Re}x\)};
    \draw[->] (0,-1.35)--(0,1.3) node[above] {\(\operatorname{Im}x\)};
    \draw[dashed] (0,0)--(58:1.7) (0,0)--(-58:1.7);
    \draw (0.45,0) arc[start angle=0,end angle=58,radius=.45];
    \node at (35:.75) {\(\theta\)};
    \draw[blue!70!black,->] (18:1.5)--(18:.72);
    \node[below] at (1.5,.12) {\(x\to0\)};
    \node[below left] at (0,0) {\(0\)};
    \node at (1,-1.55) {\(\theta=\pi/2-\eta-\varepsilon\)};
  \end{scope}
\end{tikzpicture}
\caption[Stokes 算例中的两个复平面]{Stokes 算例中的两个复平面。左图的路径绕开积分极点；右图的扇形使沿这些路径的 \(e^{-t/x}\) 衰减，且 \(e^{-1/x}\) 在 \(x\to0\) 时指数衰减。}
\label{b5:fig:D-stokes-planes}
\end{figure}
